\chapter{Your First Sprint in Logra}
\label{ch:logra}

In Part~I you drew circuits on paper and checked them in your head: you know
logic gates, multiplexers, adders, flip-flops and counters. Now you will
\emph{build} these same elements in a working environment. For an FPGA, the path from idea to
board is too long to walk by hand: you have to describe the circuit in
text, check it, see what it turned into, and only then
load it into the chip. One program does all of these steps---\term{Logra}.
It opens in a browser, needs nothing installed, and shows
every stage of the path we studied in Chapter~\ref{ch:flow}.

Open-source tools work inside Logra: Yosys turns text into a circuit,
nextpnr places it on the chip, Project Apicula packs the bitstream, and
Icarus Verilog checks the behavior. The environment only gathers them together and draws
what they computed.

\begin{buildbox}[Build it first]
Open Logra and click ``New project''. Choose ``From an example'',
find the example ``Running light'' and open it. Without changing anything, click
``Synthesize'' (Ctrl+Enter) in the top right corner. While the build runs, recall
the ring register from Part~I: six flip-flops with a single one passed
around the loop. Later in this chapter we will walk through the tabs and find each part of this circuit.
\end{buildbox}

\section{A Map of the Environment}

On the left side of Logra is a column of tabs. They are grouped in the order in
which we travel the path from idea to board: ``Learn'', ``Create'', ``Check'' and
the group ``Real board'', which is collapsed by default. Tabs with a finished result
are marked with a colored dot; the ``Pins'' tab shows on the right the number of assigned pins.

\begin{table}[H]
\centering
\caption{Logra groups and tabs, and the book chapters that cover them}
\small
\begin{tabularx}{\textwidth}{l l X l}
\toprule
\textbf{Group} & \textbf{Tab} & \textbf{What you do there} & \textbf{Chapter} \\
\midrule
Learn & Learn & Step-by-step lessons with questions and solution checking (the ``Your workshop'' course, 31 lessons) & --- \\
\midrule
Create & Code & Write a circuit description in Verilog & \ref{ch:verilog} \\
 & Designer & Build a circuit from logic gates with the mouse; Verilog writes itself & \ref{ch:logra} \\
 & Pins & Assign the circuit's ports to package pins & \ref{ch:board} \\
\midrule
Check & RTL & View the gate-level circuit built by the synthesizer & \ref{ch:bridge} \\
 & Waveforms & Apply signals and view timing diagrams & \ref{ch:sim} \\
 & Lab & Connect parts to the board and run the circuit ``as on a desk'' & \ref{ch:examples} \\
\midrule
Real board & Gowin & See LUTs, flip-flops and adder blocks after synthesis & \ref{ch:inside} \\
 & Chip & Place and route, frequency report, bitstream & \ref{ch:flow} \\
 & Connect & Load onto the board over USB, UART terminal & \ref{ch:flow} \\
\bottomrule
\end{tabularx}
\end{table}

In the top row on the right are the ``New project'' button, the language
switch (RU/EN), the theme switch, help, and the main button
\textbf{``Synthesize''} (Ctrl+Enter), which runs the whole flow. In
the bottom status bar you will see ``Yosys ready'' (this is the name of the program
Yosys, not of the environment) and ``Icarus Verilog ready''.

\begin{diybox}[A map of the environment]
\textbf{Task.} Go through all the tabs and make sure you understand the
groups.\\
\textbf{In Logra.} Open the ``Running light'' project and click ``Synthesize''.\\
\textbf{Expected result.} The RTL, Gowin and Chip tabs now have
colored dots: their results are ready; the ``Pins'' tab shows the number of assigned pins.\\
\textbf{Check.} Expand the ``Real board'' group (the arrow to the left of its
name) and find the three tabs in it. Write down which group each of the tabs
Code, Waveforms and Chip belongs to.\\
\stars{1} Repeat: find the language switch button and switch the interface back to Russian.
\end{diybox}

\section{A New Project}

Work begins with the ``New project'' button in the top row. The window
offers five ways to start: ``From scratch'' (write Verilog yourself),
``From logic elements'' (build a circuit with the mouse), ``From an example'' (the
catalog has 275 ready-made projects), ``Import a project'' (from a file) and ``Gowin project
folder''. The buttons at the top of the window are ``Create'', ``My projects'', ``Current project''
and ``Synchronization''. We will choose ``From scratch'': you only need to give the project a name and
specify the name of the top module, usually called \code{top}.

\shot{app-wizard}{The new-project window: five ways to start}{fig:y-wizard}

The ``From logic elements'' option opens the \term{Designer}: you assemble a circuit
with the mouse from a library, and Logra writes the Verilog itself. This is the same path as
on paper in Part~I, only with instant checking. The library is modeled on
the Quartus graphic editor: gates with three, four and eight
inputs, flip-flops (D, T, JK, RS), a latch, registers and counters, RAM and ROM,
multiplexers, decoders and encoders, adders, a comparator, a multiplier and
a divider, a modulo block, Gray code, and a seven-segment display decoder.
Ready-made sets are named after familiar chips of the 74 series: 74138, 74147,
74148, 7442, 7447, 74184, 7483, 7485. The ``Preview'' button lets you check
the circuit right away: click the inputs and apply a clock.

\begin{diybox}[Your first project]
\textbf{Task.} Create a project and build an ``exclusive OR'' gate in it from the
library.\\
\textbf{In Logra.} ``New project'' $\to$ ``From logic elements''; place
two inputs, an XOR element and an output.\\
\textbf{Expected result.} In ``Preview'' the output equals 1 exactly when the
inputs differ (the truth table from the chapter on gates).\\
\textbf{Check.} Open the ``Code'' tab and find the \code{assign} line with
the operator \code{\^{}}; click ``Synthesize'' and make sure that in RTL there remains
a single XOR element.\\
\stars{2} Change the circuit to get an XNOR, and compare it with the table.
\end{diybox}

\section{Code}

The ``Code'' tab is an editor with Verilog highlighting. Errors are underlined
as you type, and saving happens automatically (``Autosaved'' appears at the bottom).
Below the editor there are the ``Build log'' and ``Diagnostics''. For our walkthrough we will take the
``Running light'' project: a one in a six-bit ring register moves to the
next LED eight times a second. This is the circuit from Chapter~\ref{ch:examples},
but the ports are named the way the board names them: \code{clk}, \code{reset} and
\code{led\_n}.

\shot{app-code}{The code editor. The ``Synthesize (Ctrl+Enter)'' button in the corner runs the whole flow}{fig:y-code}

\begin{infobox}[Why \code{led\_n}]
The suffix \code{\_n} is a common convention: the signal is active \emph{low}.
The board's LEDs light when the pin is 0, so in the circuit we
feed them the inverted value of the ring.
\end{infobox}

\begin{diybox}[Your own parameter]
\textbf{Task.} Slow the running light down by half.\\
\textbf{In Logra.} The ``Code'' tab, the line \code{localparam DIV = 27\_000\_000 / 8;}.\\
\textbf{Expected result.} After replacing \code{8} with \code{4}, the light moves
four times a second.\\
\textbf{Check.} Click ``Synthesize'' and open the ``Lab'': the strip of
six LEDs should switch half as often. If ``Diagnostics'' shows a
nonzero number, read the message and find the typo.\\
\stars{2} Make the light run in the opposite direction (change the
order of bits in the concatenation).
\end{diybox}

\section{RTL: The Circuit Built by the Synthesizer}

Let us click ``Synthesize''. In under a second Yosys reads the text and builds the
circuit. The RTL tab (register-transfer level: a
description of a circuit as a set of registers and the logic between them) shows it the way
it is drawn on paper: rectangles for registers, standard symbols for
gates, circles with a plus for adders. This is exactly what we studied in
Part~I, only drawn automatically.

\shot{br-adder-rtl}{An RTL diagram, using a full adder as an example (Chapter~\ref{ch:bridge}): XOR, AND and OR gates, labels with code line numbers, and the inspector on the right}{fig:y-rtl}

Clicking an element opens the ``Inspector'' on the right: inputs, outputs, parameters and
the line of source text from which the element came. A double click
takes you to the editor. The mouse wheel zooms, and dragging pans the diagram.

\begin{ideabox}
The environment shows the circuit \emph{before} it is bound to a specific chip. This is an
``ideal'' circuit of gates and flip-flops. At the bottom of the panel it says: ``Before
mapping to the FPGA''. What it becomes on a real chip can be seen
on the Gowin tab.
\end{ideabox}

\begin{diybox}[Read the RTL]
\textbf{Task.} Find in the ``Running light'' diagram the ring of six flip-flops and
the counter's adder.\\
\textbf{In Logra.} The RTL tab after ``Synthesize''.\\
\textbf{Expected result.} You can see the registers, an addition element (a circle with a
plus), the comparison \code{cnt == DIV - 1}, and the inverters in front of the output \code{led\_n}.\\
\textbf{Check.} Click the adder: the ``Inspector'' should show the
code line \code{cnt + 1}; a double click takes you to that line in the editor.
\end{diybox}

\section{Pins: From a Port to a Pin}

The port \code{clk} in Verilog is just a name. To make it work, you have to say
which pin of the package the 27~MHz oscillator is connected to. This is what the
``Pins'' tab is for, with the sub-tabs ``Pin planner'' and ``.cst text''. On the left is a table of
your circuit's ports; on the right is the GW2AR-18 QN88 package with labels on all the pins. A hint
with three steps: choose a signal, choose a pin, click ``Assign''.
The ``Assign LEDs in order'' button places the six LEDs for you. The environment
writes the \code{IO\_LOC} lines into the file \code{board.cst}---the very same constraints
file that we studied in Section~\ref{sec:cst}.

\shot{app-pins}{The pin planner. For ports with suitable names, the pins are filled in automatically}{fig:y-pins}

\begin{diybox}[Find the pin]
\textbf{Task.} Determine which FPGA pin the \code{clk} port is connected to and
which line Logra wrote into the constraints file.\\
\textbf{In Logra.} The ``Pins'' tab: the sub-tabs ``Pin planner'' and ``.cst text''.\\
\textbf{Expected result.} \code{clk} on pin 4; the LEDs \code{led\_n[0..5]}
on pins 15--20 (Table~\ref{tab:pins}).\\
\textbf{Check.} Open ``.cst text'': the lines \code{IO\_LOC "clk" 4;} and
\code{IO\_LOC "led\_n[0]" 15;} should match the planner.\\
\stars{2} Deliberately spoil a port name in the \code{.cst} and see which
message the build gives (see the warning in Section~\ref{sec:cst}).
\end{diybox}

\section{Gowin: What the Circuit Turned Into}

The Gowin tab shows the result of the next step---mapping onto the
chip's primitives. Instead of abstract gates, here you have LUTs, DFF flip-flops,
ALU arithmetic blocks and I/O buffers. If you click a LUT, the ``Inspector''
shows a block ``LUT with 3 inputs'', the INIT configuration table, and a collapsible
``Truth table''. This is how you can plainly see the statement from
Chapter~\ref{ch:inside}: a logic gate in an FPGA is a table in memory.

\shot{br-adder-lut}{The same adder after mapping to Gowin primitives: instead of gates there are two LUTs, and on the right the inspector shows the LUT table and the INIT value}{fig:y-gowin}

\begin{diybox}[Find the LUT]
\textbf{Task.} Find one LUT in the circuit and compare its table with the function from
the code.\\
\textbf{In Logra.} Open the example ``A full adder'', click ``Synthesize'',
the Gowin tab; click a LUT.\\
\textbf{Expected result.} In the ``Inspector'' there is a 3-input table with the
INIT configuration: for the sum \code{s = a \^{} b \^{} cin}, the ones are in the rows
with an odd number of ones on the inputs.\\
\textbf{Check.} Compare the table with the truth table from Part~I and
list the rows where the output equals 1.\\
\stars{2} From the INIT value, determine which function is stored in the carry LUT.
\end{diybox}

\section{Waveforms: Behavior over Time}

The circuit is drawn, but does it work? The ``Waveforms'' tab runs the
Icarus Verilog simulator and draws timing diagrams of the signals
(Fig.~\ref{fig:y-wave}). For combinational circuits there is the scenario ``Sweep
inputs'': Logra itself applies all combinations of the input signals one after another. For
sequential circuits the environment itself generates a clock and a short reset pulse.
More about simulation is in Chapter~\ref{ch:sim}.

\shot{br-adder-wave}{All eight input combinations of the full adder: this is how you check a truth table}{fig:y-wave}

\subsection{Scenario: Events, Cycles and Signal Selection}

For circuits with memory, a single sweep is not enough: you need to apply the inputs at the right
moments. The ``Set up the scenario'' button opens a panel with four tabs.
\emph{Run} sets the clock input, the reset signal, the period and the number of cycles.
\emph{Inputs} stores constant values. On the \emph{Events} tab you can
add steps: ``in period 5 the input \code{c5} equals 1, in period 6---0 again''.
The \emph{Checks} tab lets you record expected values, and Logra
compares them with the simulation result.

If there are many signals, the ``Signals'' menu above the diagram lets you hide
the service nets and keep only the ones you need, and the ``Show the whole interval'' button
fits the entire simulation in the window. This is how the state machine and queue diagrams
in Chapter~\ref{ch:fsmdesign} were made: all the input pulses are set on the ``Events'' tab.

\begin{diybox}[Checking a truth table]
\textbf{Task.} Check the full adder by sweeping all the inputs.\\
\textbf{In Logra.} The example ``A full adder'', the ``Waveforms'' tab, ``Sweep
inputs''.\\
\textbf{Expected result.} Eight intervals: the pair \code{cout, s} equals the
binary sum \code{a + b + cin}.\\
\textbf{Check.} Pick any three intervals and compute the sum in your head;
compare it with the diagram.\\
\stars{2} For the ``Running light'', open ``Set up the scenario'' and increase the
number of cycles; find the moment when the one in the ring returned to the first
LED.
\end{diybox}

\section{Chip: Placement and the Bitstream}

The ``Place and build'' button on the ``Place and route'' panel runs
nextpnr. It decides exactly which cell of the chip to put each LUT and
each flip-flop on, and routes the connections between them. On the ``Placement map on
the chip'', colored squares show which part of the chip is occupied. Above the map are
the resource counters: for the ``Running light''---15 LUTs, 29 flip-flops, 22 ALUs and 8 I/O
pins. Next to them is the frequency report: nextpnr shows 338.87~MHz
against a required 27~MHz, a margin of about twelve times.

\shot{app-chip}{The placement map on the chip and the nextpnr result: 338.87 MHz against 27 MHz}{fig:y-chip}

After a successful build, the \code{.fs} files (bitstream), \code{JSON} and the log are available.

\begin{diybox}[Frequency margin]
\textbf{Task.} Find how many resources the ``Running light'' uses and what its
maximum frequency is.\\
\textbf{In Logra.} The ``Real board'' group, the ``Chip'' tab, ``Place and
build''.\\
\textbf{Expected result.} The counters LUT, FF, ALU, IO and a line like
``338.87 / 27.00 MHz. Frequency reached''.\\
\textbf{Check.} Divide the maximum frequency by 27 and write down the margin;
then raise the requirement in the ``Frequency, MHz'' field to 400 and rebuild: what does the report
say?\\
\stars{2} Compare the number of FFs with the number of bits in the code (22 in the counter, 1 for
\code{tick}, 6 in the ring).
\end{diybox}

\section{Lab}

Before loading a circuit onto a real board, you can run it on a
virtual desk. The Lab shows a Tang Nano 20K board on a breadboard.
At the top are the buttons ``Pause'', ``Clock'', ``Reset'', a speed selector, the board buttons
S1 and S2, and the ``Board time'' counter. You can connect parts to the board:
LEDs, displays, sensors, screens, memory. The circuit runs in the browser
continuously, and the parts behave like real ones, including the serial
protocols I$^2$C, SPI and UART (Chapter~\ref{ch:uart}).

\shot{app-lab}{The running light on the virtual board. On the left is a list of 85 parts}{fig:y-lab}

\begin{diybox}[The virtual desk]
\textbf{Task.} Put an LED on the desk and connect it to the circuit's output.\\
\textbf{In Logra.} The ``Lab'' tab: the part ``LED'' from the group
``LEDs'', a wire from a board pin.\\
\textbf{Expected result.} The LED blinks in time with the circuit; the buttons S1 and S2
change the state of the inputs.\\
\textbf{Check.} Click ``Pause'' and ``Clock'': the circuit steps one clock at a time, and
you can see when the state changes. In the parts search, find the word
``traffic light''.\\
\stars{2} Find the part ``Traffic light'' and connect it to the traffic light circuit from the examples.
\end{diybox}

\section{Connect and Learn}

The ``Connect'' tab loads the bitstream onto the board straight from the browser over
WebUSB. The switches at the top are ``Load the design'', ``UART terminal'' and ``Debugger''.
There are three steps: connect the board, choose a file, click ``Load into SRAM''. In
``UART terminal'' mode, the characters sent by the circuit over the
serial line arrive in the window, and you can send characters from it into the circuit.

\shot{app-conn}{Loading a bitstream onto the board}{fig:y-conn}

And the ``Learn'' tab is a self-contained course, ``Your workshop'', of thirty-one
lessons for those seeing an FPGA for the first time: from ``what is a bit'' to
your own circuits with the solutions checked by the simulator. There are modes ``With a mentor''
and ``Free tasks''.

\shot{app-learn}{The ``Your workshop'' course: lessons with checking}{fig:y-learn}

\begin{diybox}[The first lesson and loading]
\textbf{Task.} Complete the first lesson of the course and find the steps for loading onto the board.\\
\textbf{In Logra.} The ``Learn'' tab: the lesson ``What we have in hand''; then the
``Connect'' tab.\\
\textbf{Expected result.} The lesson is marked as completed; on ``Connect''
you see three steps and the ``Load into SRAM'' button (without a board it stays
unavailable).\\
\textbf{Check.} If you have a board, load the circuit and make sure the light
runs across the real LEDs; without a board, compare with the ``Lab''.
\end{diybox}

\section{Sprint Summary}

\begin{taskbox}[Task]
Create a project in Logra from scratch, name the module \code{top}, and enter the
description from Fig.~\ref{fig:y-code}. Click ``Synthesize'' and go in turn
through the tabs RTL, Gowin, Chip, Lab. On each one find:
how many flip-flops are in the circuit (29: six in the ring, twenty-two in the divider and one for the tick signal),
what maximum frequency nextpnr promises, and in what order the
LEDs light up.
\end{taskbox}

\begin{diybox}[The ``First path'' sprint]
\textbf{Sprint result.} A project with the module \code{top} that has passed synthesis,
placement and a run in the Lab, and your table of four numbers
(flip-flops, LUTs, frequency, pin number of \code{clk}).\\
\stars{3} \textbf{``Security Panel'' project stage:} create a project with this
name and save a first version in it (at least one sensor input and one
indicator LED). Acceptance criterion: synthesis passes without errors, in
``Pins'' all ports are assigned, and in the ``Lab'' the LED responds to button S1.
\end{diybox}
